\documentclass{article}

% ============================================================
% ICLR 2027 style
% ============================================================
\usepackage{iclr2027_conference,times}

% Optional notation commands distributed with the ICLR template.
% Keep this line as long as math_commands.tex is in the same folder.
\input{math_commands.tex}

% ============================================================
% Packages used by the manuscript
% ============================================================
\usepackage{amsmath}
\usepackage{amssymb}
\usepackage{bm}
\usepackage{mathtools}
\usepackage{booktabs}
\usepackage{graphicx}
\usepackage{multirow}
\usepackage{hyperref}
\usepackage{url}

% ============================================================
% Title and anonymous submission
% ============================================================
\title{Flow-Based Variational Inference for Group-Structured
Spike-and-Slab Neural Networks}

\author{Anonymous Authors}

% Keep commented for anonymous submission.
% \iclrfinalcopy

\begin{document}

\maketitle


% ============================================================
% Abstract
% ============================================================

\begin{abstract}

% TODO:
% 1. Group-structured spike-and-slab posterior inference.
% 2. Continuous latent activation with shared threshold.
% 3. Exact-zero bounded activations motivated by repeated structural gating.
% 4. Structure-aware spline posterior transport.
% 5. Applications to variable selection, hidden-unit sparsification,
%    and induced/independent neural connectivity.
% 6. Shallow posterior fidelity + multilayer MLP + real-world application.

\end{abstract}


% ============================================================
% 1. Introduction
% ============================================================

\section{Introduction}
\label{sec:introduction}

% TODO:
%
% Motivation:
% - DSS priors provide exact structural zeros and posterior inclusion
%   uncertainty, but inference is difficult in nonlinear models.
% - Neural-network sparsity is intrinsically structured:
%   predictors, hidden units, and edges are different selection objects.
% - Existing element-wise formulations do not provide a unified
%   representation of these structures.
%
% Main idea:
%
% one structural group
% <-> one continuous activation latent
% <-> one posterior inclusion event
%
% Contributions:
%
% 1. Group-structured latent spike-and-slab representation.
% 2. Two bounded exact-zero activations derived from repeated gating.
% 3. Structure-aware spline flows with full and lightweight attention.
% 4. Unified feature, hidden-unit, induced-edge, and singleton-edge sparsity.
% 5. Posterior-fidelity validation against MCMC in shallow neural networks,
%    followed by structured sparsification in multilayer MLPs.


% ============================================================
% 2. Method
% ============================================================

\section{Method}
\label{sec:method}


% ------------------------------------------------------------
% 2.1 Group LVR-DSS
% ------------------------------------------------------------

\subsection{Group-Structured Spike-and-Slab through Latent Activation}
\label{sec:group_lvr}

% TODO:
%
% Begin briefly from the scalar DSS prior:
%
% pi(beta_j)
% =
% (1 - omega) delta_0(beta_j)
% +
% omega pi_S(beta_j).
%
% Introduce the latent representation:
%
% beta_j
% =
% A_tau(V_j) S_lambda(U_j).
%
% Then immediately move to the main group formulation:
%
% beta_{G_g}
% =
% A_tau(V_g)
% S_{lambda_g}(U_{G_g}).
%
% One structural group shares one scalar activation latent V_g.
%
% Define:
%
% I_g = I(V_g > tau)
%
% and:
%
% PIP_g = P(V_g > tau | D).
%
% Emphasize:
% - group structure is the primary formulation;
% - ordinary scalar/element-wise DSS is the singleton-group special case;
% - the shared threshold controls global sparsity;
% - continuous slab variables retain within-group effect uncertainty.


% ------------------------------------------------------------
% 2.2 Neural structures
% ------------------------------------------------------------

\subsection{Structured Sparsity in Neural Networks}
\label{sec:nn_structure}

% TODO:
%
% Define the stacked multilayer network:
%
% h^(1) = phi(W_1 x + b_1)
%
% h^(ell) = phi(W_ell h^(ell-1) + b_ell)
%
% f_theta(x) = W_{L+1} h^(L) + b_{L+1}.
%
% Use ReLU as the ordinary hidden activation.
% ReLU here concerns function representation, not structural gating.


\subsubsection{Feature Groups}
\label{sec:feature_groups}

% TODO:
%
% Define:
%
% \tilde{x}_k = g_k^F x_k
%
% so that
%
% I_k^F = 0
%
% removes predictor x_k from the entire network.
%
% Explain this as nonlinear Bayesian variable selection.
%
% For the first layer, the equivalent parameter group corresponds to
% the complete column W_1[:, k].


\subsubsection{Hidden-Unit Groups}
\label{sec:unit_groups}

% TODO:
%
% Define unit activation:
%
% h_{\ell j}
% =
% g_{\ell j}^U
% phi(
% W_{\ell,j:} h^{(\ell-1)}
% +
% b_{\ell j}
% ).
%
% Explicitly include b_{\ell j}.
%
% In particular, b_2 is part of the second-layer unit preactivation.
%
% Main interpretation:
%
% unit selection = activation selection.
%
% Do not define units through mutually exclusive parameter ownership,
% because adjacent layers share connecting weights.


\subsubsection{Feature--Unit Induced Connectivity}
\label{sec:induced_edges}

% TODO:
%
% First layer:
%
% I_{1,jk}^{E,ind}
% =
% I_k^F I_{1j}^U.
%
% Hidden-to-hidden:
%
% I_{\ell,ji}^{E,ind}
% =
% I_{\ell-1,i}^U I_{\ell,j}^U.
%
% Output:
%
% I_{out,i}^{E,ind}
% =
% I_{L,i}^U.
%
% Important:
% induced edge activity is a posterior-derived connectivity quantity,
% not an additional gate multiplied again into the forward model.
%
% This gives:
%
% feature sparsity
% +
% unit sparsity
% =>
% induced edge/path sparsity.


\subsubsection{Independent Edge Groups}
\label{sec:edge_groups}

% TODO:
%
% Treat each neural-network weight as a singleton group:
%
% G_{\ell,ji}^E
% =
% {W_{\ell,ji}}.
%
% Define:
%
% \widetilde W_{\ell,ji}
% =
% g_{\ell,ji}^E W_{\ell,ji}.
%
% This recovers ordinary edge-wise neural sparsification as the
% singleton-group special case of the general group formulation.
%
% Bias terms are not edges.


% ------------------------------------------------------------
% 2.3 Exact-zero gates
% ------------------------------------------------------------

\subsection{Exact-Zero Activations from Repeated Structural Gating}
\label{sec:exact_zero_activation}

% TODO:
%
% This section should derive the two bounded gate constructions
% from the same hard structural indicator.
%
% Let
%
% m = V - tau,
%
% I = I(m > 0).


\subsubsection{Repeated Hard Gating and Indicator Idempotence}
\label{sec:indicator_idempotence}

% TODO:
%
% For a hidden unit, the same structural indicator may appear on
% both incoming and outgoing contributions:
%
% I_j W_{2j}
% ReLU[
% I_j (W_{1j}^T x + b_j)
% ].
%
% By positive homogeneity of ReLU:
%
% =
% I_j^2 W_{2j}
% ReLU(W_{1j}^T x + b_j).
%
% For a hard indicator:
%
% I_j^2 = I_j.
%
% This provides two different routes to a differentiable,
% exact-zero representation.


\subsubsection{Bounded RePU from Continuous Gate Multiplication}
\label{sec:normalized_requ}

% TODO:
%
% Route 1:
% relax the repeated hard indicator before applying idempotence.
%
% Replace
%
% I_j -> (m_j)_+.
%
% Then:
%
% I_j^2
% ->
% (m_j)_+^2,
%
% giving a ReQU/RePU-type gate.
%
% Use the bounded default:
%
% G_N(m; rho)
% =
% (m_+)^2 /
% [rho^2 + (m_+)^2].
%
% Main experiments:
%
% rho = 1.
%
% Properties:
% - exact zero for m <= 0;
% - bounded in [0,1);
% - continuously adjustable on the active side.
%
% This is the default activation used in the main experiments.


\subsubsection{Smooth Exact-Zero Gate after Indicator Idempotence}
\label{sec:smooth_step}

% TODO:
%
% Route 2:
% first use the hard identity
%
% I^2 = I,
%
% and then smoothly approximate the resulting single indicator.
%
% Define:
%
% G_S(m; delta)
% =
% 0,                         m <= 0
%
% exp(-delta/m)
% -----------------------------------,   0 < m < delta
% exp(-delta/m) + exp[-delta/(delta-m)]
%
% 1,                         m >= delta.
%
% Main experiments:
%
% delta = 1.
%
% Properties:
% - exact zero;
% - exact one;
% - C^\infty transition;
% - same hard structural event I(m > 0).
%
% Conceptual distinction:
%
% normalized ReQU:
% relax first, then multiply;
%
% smooth-step:
% simplify I^2 = I first, then relax.


% ------------------------------------------------------------
% 2.4 Spline flow
% ------------------------------------------------------------

\subsection{Structure-Aware Spline Posterior Transport}
\label{sec:spline_flow}

% TODO:
%
% Collect latent variables:
%
% xi = (U, V, tau)
%
% or the corresponding concatenated feature/unit/edge latent blocks.
%
% Base distribution:
%
% z_0 ~ N(mu_0, diag(sigma_0^2)).
%
% Flow:
%
% xi = T_psi(z_0).
%
% Use complementary coupling masks:
%
% xi_A' = xi_A
%
% xi_j'
% =
% T_{theta_j(xi_A)}(xi_j),
% j in B,
%
% where T is a monotone rational-quadratic spline.
%
% The conditioner determines theta_j:
% spline widths, heights, and derivatives.
%
% Two conditioner variants are considered below.


\subsubsection{Full-Attention Spline Flow}
\label{sec:full_attention_spline}

% TODO:
%
% High-capacity reference conditioner.
%
% Fixed-side coordinates contain:
% - current latent value;
% - coordinate identity;
% - latent role;
% - group identity;
% - layer / unit / feature / edge metadata when applicable.
%
% Structure:
%
% fixed value + metadata tokens
% ->
% fixed-side self-attention
% ->
% target cross-attention
% ->
% spline parameter readout.
%
% Target queries may use structural metadata but must not contain
% the current target random value xi_j.
%
% Therefore spline parameters remain functions of xi_A only,
% preserving the triangular coupling Jacobian.


\subsubsection{Lightweight Attention Spline Flow}
\label{sec:light_attention_spline}

% TODO:
%
% Efficient conditioner without self-attention or FFN.
%
% Metadata embedding:
%
% e_i.
%
% For head h, target query:
%
% q_j^(h)
% =
% Q^(h) e_j.
%
% Value-dependent source key:
%
% k_i^(h)
% =
% K^(h)
% [
%   e_i
%   xi_i
% ].
%
% Attention:
%
% alpha_{ji}^{(h)}
% =
% softmax_{i in A}
% [
% (q_j^(h))^T k_i^(h) / sqrt(d_h)
% ].
%
% Deterministic source value:
%
% v_i
% =
% [
%   xi_i
%   xi_i^2
%   P_v e_i
% ].
%
% Context:
%
% c_j^(h)
% =
% sum_{i in A}
% alpha_{ji}^{(h)} v_i.
%
% Use a small number of heads, e.g. H_attn = 2.
%
% Linear parameter heads map the resulting contexts to:
% - spline width logits;
% - spline height logits;
% - spline derivative logits.
%
% No self-attention.
% No Transformer FFN.
% No MLP conditioner.
%
% Interpretation:
%
% attention determines which fixed latent coordinates matter;
% spline transport determines how the target coordinate is transformed.


% ------------------------------------------------------------
% 2.5 VI
% ------------------------------------------------------------

\subsection{Variational Inference and Representation Warmup}
\label{sec:vi}

% TODO:
%
% Variational density:
%
% q_psi(xi)
%
% induced by the spline flow.
%
% Optimize:
%
% L(psi)
% =
% E_{q_psi(xi)}
% [
% log p(D | xi)
% +
% log p(xi)
% -
% log q_psi(xi)
% ].
%
% Include exact spline log-Jacobian terms through
% change of variables.


\subsubsection{Representation Warmup}
\label{sec:warmup}

% TODO:
%
% During the first E_w epochs:
%
% g_g^{like} = 1
%
% for all selectable structures in the likelihood.
%
% Important:
% - posterior latent samples are not modified;
% - prior is evaluated on the actual latent draw;
% - log q is evaluated on the actual latent draw;
% - only the likelihood temporarily sees all structural gates as active.
%
% After warmup, activate the learned gate normally.
%
% Explain this as separating representation learning from
% subsequent structural sparsification.


% ------------------------------------------------------------
% 2.6 PIP
% ------------------------------------------------------------

\subsection{Posterior Structural Inference and Inclusion Probabilities}
\label{sec:pip}

% TODO:
%
% Posterior samples:
%
% xi^(r) = T_{\hat psi}(z_0^(r)).
%
% Regardless of continuous optimization gate,
% recover the hard structural state using:
%
% I_g^(r)
% =
% I(V_g^(r) > tau^(r)).
%
% Group PIP:
%
% \widehat{PIP}_g
% =
% (1/R) sum_r I_g^(r).


\subsubsection{Feature Inclusion Probabilities}

% TODO:
%
% PIP_k^F
% =
% P(I_k^F = 1 | D).


\subsubsection{Hidden-Unit Inclusion Probabilities}

% TODO:
%
% PIP_{\ell j}^U
% =
% P(I_{\ell j}^U = 1 | D).


\subsubsection{Induced Edge Inclusion Probabilities}

% TODO:
%
% First layer:
%
% PIP_{1,jk}^{E,ind}
% =
% E[
% I_k^F I_{1j}^U
% | D
% ].
%
% Hidden layers:
%
% PIP_{\ell,ji}^{E,ind}
% =
% E[
% I_{\ell-1,i}^U I_{\ell,j}^U
% | D
% ].


\subsubsection{Independent Edge Inclusion Probabilities}

% TODO:
%
% PIP_{\ell,ji}^E
% =
% P(I_{\ell,ji}^E = 1 | D).


% ============================================================
% 3. Experiments
% ============================================================

\section{Experiments}
\label{sec:experiments}


% ------------------------------------------------------------
% 3.1 Shallow neural network
% ------------------------------------------------------------

\subsection{Posterior and Selection Recovery in Shallow Neural Networks}
\label{sec:shallow}

% Main purpose:
%
% posterior fidelity against matched MCMC in a setting where
% reliable posterior reference computation remains feasible.


\subsubsection{Simulation Design}
\label{sec:shallow_setup}

% TODO:
%
% Main setting:
%
% n = 240
% p = 6
% H_fit = 5
% H_true = 2
% p_true = 2
%
% active features:
%
% x_0, x_3.
%
% Gaussian response with sigma^2 = 1.
%
% Keep nonlinear teacher and train/validation/test splitting.
%
% Flow:
% K = 4
% warmup = 500
% total epochs = 2000
% R_train = 100.
%
% Matched MCMC is used as the posterior reference.


\subsubsection{Variable Selection}
\label{sec:shallow_feature}

% TODO:
%
% Evaluate feature-group posterior.
%
% Report:
% - feature PIPs;
% - VI vs MCMC;
% - truth active variables x_0 and x_3;
% - active posterior density;
% - function MSE.
%
% Feature identity is identifiable, so support recovery may be
% formally assessed here.


\subsubsection{Hidden-Unit Selection}
\label{sec:shallow_unit}

% TODO:
%
% Evaluate hidden-unit groups.
%
% Report:
% - ranked VI PIPs;
% - ranked MCMC PIPs;
% - posterior density of active functional components;
% - function recovery.
%
% Explicitly state that teacher-unit identities represent a
% canonical generating representation rather than identifiable
% posterior labels.


\subsubsection{Posterior and Function Recovery}
\label{sec:shallow_recovery}

% TODO:
%
% This is the main posterior-fidelity subsection.
%
% Metrics:
%
% D_SKL,A
% D_JS,0
% D_SKL,joint
%
% plus signal MSE / R^2.
%
% Figures:
% - marginal active posterior;
% - joint 2D HDR posterior;
% - PIP comparison;
% - posterior function comparison.
%
% MCMC diagnostics are reported in the appendix.


\subsubsection{Activation and Spline-Flow Ablations}
\label{sec:shallow_ablation}

% TODO:
%
% Main comparisons:
%
% Gate:
% - normalized ReQU;
% - smooth-step.
%
% Conditioner:
% - full-attention spline;
% - lightweight-attention spline.
%
% If retained:
% affine / ReLU results go primarily to the appendix.


% ------------------------------------------------------------
% 3.2 Multilayer MLP
% ------------------------------------------------------------

\subsection{Structured Selection in Multilayer MLPs}
\label{sec:mlp}

% Main model:
%
% stacked + feature-unit-induced-edge
%
% Main objective:
%
% structured sparsification and function recovery,
% rather than exact hidden-neuron label recovery.


\subsubsection{Multilayer Architecture}
\label{sec:mlp_architecture}

% TODO:
%
% Main stacked MLP:
%
% p
% -> H_1
% -> H_2
% -> 1.
%
% Example default:
%
% hidden_dims = (5,5).
%
% Deeper stress test:
%
% hidden_dims = (6,5,4).
%
% Main model:
% stacked architecture.
%
% Embedding-output architecture is retained as an architectural
% ablation rather than the default model.


\subsubsection{Variable Selection}
\label{sec:mlp_feature}

% TODO:
%
% Feature-group selection in nonlinear multilayer MLP.
%
% Report:
% - original-variable PIPs;
% - recovery of active x_0 and x_3;
% - signal MSE;
% - structure stability across repeated VI seeds.


\subsubsection{Feature--Unit Induced Connectivity}
\label{sec:mlp_induced}

% TODO:
%
% Main structured sparse MLP.
%
% Report:
% - feature PIPs;
% - layer-wise unit PIPs;
% - expected active units;
% - induced edge PIPs;
% - induced edge density;
% - network connectivity visualization.
%
% Main functional result:
%
% signal MSE versus induced edge density.


\subsubsection{Independent Edge Sparsification}
\label{sec:mlp_edge}

% TODO:
%
% Each edge is treated as a singleton group.
%
% Report:
% - edge PIP matrices;
% - expected edge density;
% - signal MSE;
% - comparison against induced connectivity.
%
% Do not use deep edge-label truth as the primary metric.


\subsubsection{MSE--Sparsity Trade-Off}
\label{sec:mse_sparsity}

% TODO:
%
% For posterior inclusion thresholds c:
%
% construct sparse posterior structures
%
% and plot:
%
% signal MSE
% versus
% retained edge density.
%
% Mark the median probability model:
%
% c = 0.5.
%
% Compare:
% - feature-unit induced connectivity;
% - independent edge sparsification.


\subsubsection{Stacked versus Embedding--Output Parameterizations}
\label{sec:embedding_ablation}

% TODO:
%
% Present the architectural ablation.
%
% Main current conclusion:
%
% - feature-unit induced grouping substantially alleviates
%   collapse in the redundant embedding-output parameterization;
% - however, stacked induced remains substantially better in
%   function recovery and computational efficiency.
%
% Therefore stacked + induced is the primary MLP architecture.


\subsubsection{Network Depth}
\label{sec:depth}

% TODO:
%
% Test whether the same group activation and spline-flow construction
% remains effective when moving from two hidden layers to a modestly
% deeper MLP.
%
% This is a scalability/generalization check,
% not a search for the optimal depth.


% ------------------------------------------------------------
% 3.3 Real-world longitudinal data
% ------------------------------------------------------------

\subsection{Real-World Longitudinal Data Application}
\label{sec:real_data}


\subsubsection{Dataset and Prediction Task}
\label{sec:real_dataset}

% TODO:
%
% Select a longitudinal dataset with:
% - repeated observations per subject;
% - interpretable longitudinal covariates;
% - a regression or otherwise compatible predictive outcome;
% - natural grouping of measurements or variable domains.


\subsubsection{Longitudinal Group Construction}
\label{sec:real_groups}

% TODO:
%
% Define meaningful groups from:
% - repeated measurements of the same biomarker;
% - temporal measurements;
% - clinical domains;
% - variable families.
%
% Avoid reducing the real-data demonstration to arbitrary
% singleton-variable sparsification unless needed as a baseline.


\subsubsection{Sparse Multilayer Posterior}
\label{sec:real_sparse_mlp}

% TODO:
%
% Main model:
%
% stacked + feature-unit-induced-edge
%
% with:
%
% normalized ReQU
% +
% spline posterior transport.
%
% Report:
% - group PIPs;
% - unit sparsity;
% - induced connectivity;
% - posterior predictive function.


\subsubsection{Predictive and Structural Results}
\label{sec:real_results}

% TODO:
%
% No support truth exists in real data.
%
% Therefore emphasize:
% - held-out prediction;
% - selected longitudinal groups;
% - sparse network structure;
% - posterior inclusion uncertainty;
% - interpretability of selected longitudinal effects.


% ============================================================
% 4. Conclusion
% ============================================================

\section{Conclusion}
\label{sec:conclusion}

% TODO:
%
% Summarize the unified framework:
%
% group-structured latent spike-and-slab
% +
% bounded exact-zero activation
% +
% structure-aware spline posterior transport.
%
% Emphasize the progression:
%
% variable selection
% ->
% hidden-unit sparsification
% ->
% induced neural connectivity
% ->
% multilayer sparse neural networks.
%
% Do not describe grouped or nonlinear modeling as future work;
% they are now central contributions of the paper.


% ============================================================
% Required / recommended ICLR statements
% ============================================================

\subsection*{AI use statement}

% TODO: replace with the final ICLR-compliant disclosure.
We used generative AI tools to assist with manuscript editing and
software development. All AI-assisted text, mathematical expressions,
and code were reviewed and verified by the authors. The authors take
full responsibility for the final content of this work.


\subsection*{Reproducibility statement}

% TODO:
% Summarize:
% - simulation generators;
% - model and flow implementation;
% - posterior sampling;
% - MCMC reference settings for shallow experiments;
% - seeds and computing environments;
% - supplementary code and data.


% ============================================================
% References
% ============================================================

\bibliography{references}
\bibliographystyle{iclr2027_conference}


% ============================================================
% Appendix
% ============================================================

\appendix


\section{Additional Methodological Details}
\label{app:method}

% TODO:
%
% Potential material:
% - complete derivation of normalized ReQU and smooth-step;
% - smooth-step C^\infty properties;
% - exact spline formulas;
% - coupling masks;
% - metadata definitions;
% - full-attention architecture details;
% - lightweight-attention implementation details;
% - embedding-output parameterization.


\section{Additional Shallow-Network Experiments}
\label{app:shallow}

% TODO:
%
% - full MCMC diagnostics;
% - additional posterior-density plots;
% - affine coupling ablations;
% - ReLU / unbounded gate controls if retained;
% - additional seeds.


\section{Additional Multilayer MLP Experiments}
\label{app:mlp}

% TODO:
%
% - full embedding-output results;
% - independent edge PIP matrices;
% - induced connectivity matrices;
% - additional depths;
% - additional VI seeds;
% - runtime and parameter-count tables.


\section{Real-World Data Details}
\label{app:real}

% TODO:
%
% - preprocessing;
% - missing-data handling;
% - longitudinal group definitions;
% - train/validation/test split;
% - full selected-group tables.


\end{document}